\section{Random domains and predictable endpoint bounds}\label{sec:noise-extensions}

The fixed-list proof is complete. An independent random choice of observed
labels gives concentration in coordinate subcubes much smaller than a
half-cube. A second refinement builds endpoint caps on one common event,
giving simultaneous pointwise control of the conditional probabilities.

\subsection{Random blockers and disjoint subcubes}

When labels are sampled uniformly, every unconditioned available set is a
uniform subset of its prescribed size. This allows concentration in
subcubes much smaller than half the cube. It yields a different route
to averaged marginal bounds.

\begin{lemma}[Decorrelation across preceding-coordinate subcubes]
\label{noise:subcubes}
Suppose the blockers are a uniform subset of starting labels, chosen
independently of the shuffle, and $m$ labels remain unblocked. Fix
$f_t$ to be the centered conditional vector of a tagged unblocked card,
so $\|f_t\|_2^2\le1$, and let $1\le L<d$.
At time $s\ge L$, partition $V$ into subcubes varying the
$L$ directions used by layers $s-L,\ldots,s-1$. If, with probability
at least $1-\eta$, each such subcube at time $s-L$ has a fraction
at least $r$ of available positions, then
\[
 w_s\le(1-r)a_s+\eta.
\]
\end{lemma}

\begin{proof}
Condition on the initial label choice and all switch history through
$s-L$. Each of these subcubes now evolves with its own independent
coins. The update rule for $f$ is local, so $f_s(x)^2$ depends on the
coins in the subcube of $x$. Its partner $x^{i_s}$ lies in a different
subcube, since the current coordinate has not been used among the
preceding $L$ coordinates. The indicator that the partner is available is conditionally
independent of $f_s(x)^2$. Its conditional mean is the old density of
available positions in its subcube: averaging in each of that subcube's $L$ coordinates
makes the expected occupancy constant. On the asserted event this
gives the factor $1-r$. On its complement the sum of squared masses
is at most one. Averaging proves the result. The argument uses
unconditional concentration of the old available set, rather than claiming
that it is uniform after path conditioning.
\end{proof}

\begin{proposition}[Averaged complements with arbitrary fixed density]
\label{noise:averaged-general}
Let $b\ge2$ be a fixed power of two, $q=1/b$, $T=100bd$, and
$A\subseteq V$ be a uniform subset of size $n-n/b$. If $\mu$ is the
time-$T$ permutation law, then
\begin{equation}
 \E_A\left\|\mathcal L_{g\sim\mu}(g|_A)
             -\Unif(\operatorname{inj}(A,V))\right\|_{\TV}=o(1).
 \label{noise:averaged-injections}
\end{equation}
\end{proposition}

\begin{proof}
Choose a uniform ordered list with underlying set $A$. At any exposure
index, the earlier labels are blockers and the next label is tagged;
the number of available positions is at least $qn$. Put $L=\lfloor d/2\rfloor$.
At every fixed time the available set is unconditionally uniform, since a
shuffle permutation independent of the initial selection preserves
uniform subset sampling. Its chance of having fraction below $q/2$
in any one of the preceding-coordinate subcubes is at most
\begin{equation}
 \eta_d=n(n+1)e^{-q2^L/8}.
 \label{noise:binomial-conditioned-tail}
\end{equation}
Indeed sample independent Bernoulli indicators of parameter $m/n$,
then condition their total to equal $m$. This conditioning gives a
uniform $m$-subset, and the conditioning event has probability at
least $1/(n+1)$: the value $m$ is a mode of the binomial distribution
with parameter $m/n$. The degenerate parameter one is immediate.
Before conditioning, the lower-tail Chernoff bound for a subcube
of size $2^L$ is at most $e^{-q2^L/8}$. A union bound over at most
$n$ subcubes proves \eqref{noise:binomial-conditioned-tail}.

Lemma~\ref{noise:subcubes} gives $w_s\le\rho a_s+\eta_d$ for
$s\ge L$, with $\rho=1-q/2$. Set $t_0=d+L$ and $\theta=1-q/8$.
Since $\rho/(2\theta-1)<1$, the comparison induction gives
\[
 a_t\le\theta^{t-t_0}+2\eta_d/q.
\]
At $T=100bd$ the geometric term is at most
$\exp(-(100b-1.5)qd/8)$, and hence $a_T=o(n^{-5})$ uniformly in
the exposure index. The sequential comparison used in
Proposition~\ref{noise:uniform-marginals}, now averaged over the
initial list, bounds its mean tuple distance by
$n^{3/2}\sqrt{a_T}/2=o(1)$. The order of the chosen starting labels
does not affect the distance for their restriction map, proving
\eqref{noise:averaged-injections}.
\end{proof}

The preceding argument already includes a random list of $3n/4$ labels.
For this density, a direct exponential moment gives another tail bound
without introducing independent Bernoulli variables and conditioning on
their total. We record that calculation with the energy rate used in the
four-subgroup application below.

\begin{proposition}[A direct moment bound for three quarters of the labels]
\label{noise:covariance-marginal}
At $T=1000d$, the mean total-variation distance for the images of a
uniform ordered list of $3n/4$ starting labels from a uniform injective
tuple tends to zero. In its conditional one-card construction,
\begin{equation}
 w_s\le\frac78a_s+\eta_d\quad(s\ge L),\qquad
 \eta_d=n\left((7/4)2^{-7/8}\right)^{2^L},\quad
 L=\lfloor d/2\rfloor,
 \label{noise:covariance-tail}
\end{equation}
and $a_T=o(n^{-10})$.
\end{proposition}

\begin{proof}
The density of available positions is at least $1/4$. If $B_C$ counts blockers in a
fixed preceding-coordinate subcube $C$, uniform sampling gives
\[
 \E 2^{B_C}
 =\E\prod_{x\in C}(1+\one_{\{x\text{ blocked}\}})
 \le\sum_{J\subseteq C}(3/4)^{|J|}=(7/4)^{|C|}.
\]
Markov's inequality therefore bounds the probability that $B_C$
exceeds $7|C|/8$ by $((7/4)2^{-7/8})^{|C|}$. Union over at most
$n$ subcubes and Lemma~\ref{noise:subcubes} prove
\eqref{noise:covariance-tail}. The base is strictly less than one.
Using the exact covariance recursion \eqref{noise:covariance}, or its
trace identity \eqref{noise:memory-equality}, gives
\[
 a_t\le(31/32)^{t-2d}+8\eta_d.
\]
The comparison is valid because $(7/8)/(2(31/32)-1)=14/15<1$
and $(7/8)8+1=8$. At $T=1000d$ this is $o(n^{-10})$.
The expected conditional variation per list entry is at most
$\sqrt{ma_T}/2\le\sqrt{na_T}/2$. Summing as before proves the
claimed mean tuple bound.
\end{proof}

\subsection{Predictable endpoint bounds}

Small expected variation can be turned into bounded coset densities by
trimming. A second construction produces these bounds directly on one
common high-probability event. The event for the next card must be
measurable before that card's path is exposed.

\begin{proposition}[A common subprobability with predictable coset bounds]
\label{noise:predictable-cosets}
For all sufficiently large $d$, let $T=2048d$ and let $\mu$ be its
permutation law. There is a
partition $D_1,\ldots,D_8$ of the starting positions into sets of size
$n/8$ and a subprobability $0\le f\le\mu$ such that
\begin{equation}
 \sum_g f(g)\ge1-8n^{-11},\qquad
 f(gH_i)\le2\frac{|H_i|}{|G|}
 \quad(g\in G,\ 1\le i\le8),
 \label{noise:predictable-cap}
\end{equation}
where $H_i=\Sym(D_i)$ fixes the complement pointwise.
\end{proposition}

\begin{proof}
Choose a uniform ordered partition and independently uniform orders
of its eight complements. For one of these random orders and a fixed exposure index, let
earlier labels be blockers, put $p=1/8$, and use the centered vector
$f_t$ already defined. Unconditionally its available set is uniform and
has size at least $pn$.

Write $s=\ell d+j$, with $\lceil d/2\rceil\le j\le d-1$.
Conditional on the history at $\ell d$, the first $j$ coordinate
updates evolve independently in their $j$-dimensional subcubes.
The argument of Lemma~\ref{noise:subcubes}, with this pass boundary
in place of $s-L$, gives $w_s\le(1-p/2)a_s+\delta_n$, where
\[
 \delta_n=n\exp(-p^2\sqrt n/16).
\]
For completeness, in a subcube of size $r$ let $B$ count blockers.
For every $\vartheta>0$, expansion over subsets and sampling without
replacement give
\[
 \E e^{\vartheta B}
 \le[1+(1-p)(e^{\vartheta}-1)]^r.
\]
Take $\vartheta=p/4$ and use
$e^\vartheta-1\le\vartheta+\vartheta^2$. The logarithm of the
Markov bound for $B>(1-p/2)r$ is at most
\[
 r\{-\vartheta(1-p/2)+(1-p)(\vartheta+\vartheta^2)\}
 \le-rp^2/16.
\]
Here $r=2^j\ge\sqrt n$, and a union bound proves the displayed
$\delta_n$. At all other times simply use $w_s\le a_s$.

Starting with $a_d\le1$, the one-step memory inequality
\eqref{noise:one-step-memory} now contracts by $1-p/4$ at each
late step, and otherwise does not increase except for the additive
bound. There are $(2048-1)\lfloor d/2\rfloor$ late steps after
time $d$. Thus
\begin{equation}
 a_T\le(1-p/4)^{2047\lfloor d/2\rfloor}+T\delta_n
 \le n^{-18}
 \label{noise:predictable-energy}
\end{equation}
eventually. For the last inequality one may use
$\lfloor d/2\rfloor\ge d/3$ and $p/4=1/32$.

For each query define $E_a=\{\|f_T\|_2\le n^{-3}\}$.
Its failure probability is at most $n^{-12}$ by
\eqref{noise:predictable-energy}. There are at most $8n$ queries,
so their simultaneous success has probability at least $1-8n^{-11}$,
averaging over the initial choices and the shuffle. Fix choices of
partition and orders attaining at least this probability. Let $f(g)$
be the probability of permutation $g$ together with simultaneous
success. This proves its domination and mass assertion.

It remains to prove its pointwise coset bounds. In one fixed order,
let $\mathcal B_a$ reveal the full paths through $T$ of the first
$a$ labels. The event $E_a$ for query $a$ depends only on
$\mathcal B_{a-1}$, since the initial tag is fixed and its conditional
law is computed from earlier paths. On this event the conditional
probability of a specified next endpoint is at most
\[
 \frac1{n-a+1}+n^{-3}
 \le\frac{1+n^{-2}}{n-a+1}.
\]
For specified distinct endpoints, let $S_a$ be the event that the
$a$th card reaches its specified endpoint. In the product of
$\one_{S_a}\one_{E_a}$, condition first on
$\mathcal B_{n-n/8-1}$ and remove the last factor using this bound;
then continue backwards. All previous factors are measurable at
the conditioning step. This proves that the probability of the
specified endpoint list together with all its good events is at most
\[
 (1+n^{-2})^{n-n/8}\frac{(n/8)!}{n!}
 \le2\frac{|H_i|}{|G|}.
\]
The subprobability $f$ also requires the other orders' good events,
so it is bounded by this probability. A coset $gH_i$ specifies
exactly this complement endpoint list, proving
\eqref{noise:predictable-cap}.
\end{proof}


\subsection{The four corresponding full-permutation applications}

The information just obtained has three different forms: a uniform
list estimate, an average over lists, and one subprobability with
predictable pointwise caps. They enter distinct Fourier arguments.
The statements below preserve the quantitative conclusions of those
arguments; all times count physical shuffles.

\begin{theorem}\label{noise:mixing}
For the Thorp shuffle on $n=2^d$, each of the following methods proves
worst-start total-variation convergence to uniform at the indicated time:
\begin{enumerate}
\item the sixteen-subgroup orbit-span argument at $8192d$;
\item predictable coset caps at $16384d$;
\item diagram peeling at $400\cdot2^{22}d$;
\item the covariance and four-subgroup coefficient argument at $16000d$.
\end{enumerate}
The $2048d$ fixed-list and isotypic route was completed in
Section~\ref{sec:first-route}.
\end{theorem}

\begin{proof}
Write $C_u=\sup_{m\ge1}\sum_{\tau\vdash m}D_\tau^{-u}$ for the
finite constants proved in~\cite[Theorem~\ref*{l-l:degree-all-powers}]{partialFourier}.
For the separate peeling argument put $C_{\rm peel}=512^2$.
As before, every untrimmed positive-time shuffle block has zero sign
expectation. All the lifting statements invoked below are proved
independently in~\cite{partialFourier}.

For the $1024d$ law, apply the first part of
Proposition~\ref{noise:uniform-marginals} with $b=16$ blocks.
Use the sequential trimming operation in
Lemma~\ref*{l-l:trim} of~\cite{partialFourier}: reduce each coset
mass exceeding its uniform mass to that mass. Each reduction preserves
the earlier caps, and the total mass removed is at most
$\varepsilon=bn^{-3/2}$. The resulting subprobability $f\le\mu$
has mass $z\ge1-\varepsilon$ and caps one. Thus $\nu=f/z$
satisfies $\|\nu-\mu\|_{\TV}\le\varepsilon$ and has cap two
when $\varepsilon\le1/2$.
Proposition~\ref*{l-l:orbit-span} of~\cite{partialFourier} bounds the operator norm by
$A D^{-3/8}$, $A=\sqrt{2bC_2}$. Since
$\|B^8\|_{\HS}\le\sqrt D\|B\|_{\op}^8$, the nonsign
sum is at most $A^{16}\sum_{D>1}D^{2-6}=o(1)$.
Its sign coefficient has absolute value at most $2bn^{-3/2}$.
The same probability comparison proves mixing at $8\cdot1024d$.

For the predictable construction, let $f\le\mu$ be supplied by
Proposition~\ref{noise:predictable-cosets}, and put $z=\sum f$.
Proposition~\ref*{l-l:isotypic} of~\cite{partialFourier}, with $B=2,b=8$, gives
$\|\widehat f(\rho)\|_{\HS}^2\le16C_2D^{-1/2}$.
Eight powers again leave an inverse-cube dimension sum. Its sign
coefficient is bounded by $1-z=o(1)$ and its trivial coefficient
is $z$. More explicitly, for any fixed $R$,
\begin{equation}
 \|f^{*R}-z^R U_G\|_1^2
 \le\sum_{\rho\ne\mathrm{triv}}
 D_\rho\|\widehat f(\rho)^R\|_{\HS}^2.
 \label{noise:subprob-plancherel}
\end{equation}
Thus the right side tends to zero for $R=8$.
The differences $\mu^{*R}-f^{*R}$ and $U_G-z^RU_G$ have
$\ell^1$ norms $1-z^R=o(1)$, since the former is nonnegative.
This proves mixing after $8\cdot2048d$.

For peeling, take $b=2^{22}$ and $T=100bd$ in
Proposition~\ref{noise:averaged-general}. A uniform equal-size
partition has uniformly distributed complements, so averaging
selects one deterministic partition whose sum of complement
errors is $o(1)$. Heavy-coset deletion gives a subprobability
$\eta\le\mu$ of mass $1-o(1)$ and cap two. Since
$b\ge4(C_{\rm peel}+2)$, Corollary~\ref*{l-l:peeling-lift} of~\cite{partialFourier} implies
\[
 \sum_{D_\rho>1}D_\rho\|\widehat\eta(\rho)^4\|_{\HS}^2
 \le(6b)^4\sum_{D_\rho>1}
 D_\rho^{1-4(1-(C_{\rm peel}+2)/b)}=o(1),
\]
because the exponent is at most $-2$ and
Theorem~\ref*{l-l:degree-all-powers} of~\cite{partialFourier} applies. The trivial and
sign coefficients, and restoration of missing mass, are handled
exactly by \eqref{noise:subprob-plancherel}. The total physical
time is $4T=400\cdot2^{22}d$.

Finally, Proposition~\ref{noise:covariance-marginal} selects,
by the same averaging argument, a deterministic partition into
four blocks with complement errors summing to $o(1)$ for the
$1000d$ law. Trim to a subprobability $\eta$ of mass $1-o(1)$
and cap two. Proposition~\ref*{l-l:four-coefficient} of~\cite{partialFourier} gives
$\|\widehat\eta(\rho)\|_{\op}\le C D^{-1/8}$.
After sixteen factors the nontrivial nonsign sum is at most
\[
 C^{32}\sum_{\rho\ne\mathrm{triv},\sgn}D^{2-16/4}
 =C^{32}\sum_{\rho\ne\mathrm{triv},\sgn}D^{-2}=o(1).
\]
Again parity and missing mass are controlled as above. This
proves the covariance route at $16\cdot1000d=16000d$.


All blocks contain a whole number of sweeps. The moving frame therefore
returns to the physical coordinates, and independent blocks have the
convolution laws used above. Equation~\eqref{eq:worst-state} supplies
worst-start uniformity in every case. The lower obstruction is
\eqref{eq:support}; it is independent of which upper-bound method is used.
\end{proof}
